\documentclass[11pt,a4paper]{article}

\usepackage[utf8]{inputenc}
\usepackage[T1]{fontenc}
% english is the main language; italian is available for the extended sommario.
\usepackage[english]{babel}
\usepackage{amsmath,amssymb,amsfonts,amsthm}
\usepackage{geometry}
\usepackage{booktabs}
\usepackage{array}
\usepackage{longtable}
\usepackage{xcolor}
\usepackage{listings}
\usepackage{hyperref}

\geometry{margin=2.5cm}
\hypersetup{colorlinks=true, linkcolor=black, urlcolor=blue, citecolor=black}

\lstdefinestyle{code}{
  basicstyle=\ttfamily\small,
  breaklines=true,
  frame=single,
  columns=fullflexible,
  backgroundcolor=\color{gray!5}
}
\lstset{style=code}

\newtheorem{proposition}{Proposition}
\newtheorem{definition}{Definition}

\title{SemaGraph\\[2pt]
\large A deterministic, cache-resident kernel for classifying and compressing
state-transition trajectories over a six-bit alphabet}
\author{Technical report --- version 0.7 (\texttt{semagraph-64})}
\date{June 2026}

\begin{document}
\maketitle

\begin{abstract}
SemaGraph is a deterministic kernel that represents observed change as finite
trajectories over a small fixed alphabet and classifies and compresses those
trajectories using only native integer operations. Each state is a six-bit word
($Q_6$, $64$ states); each direct transition is an XOR mutation mask; the complete
transition basis is the $64\times64=4096$ ordered pairs. Per-transition metadata
(Hamming distance, per-module distances, module scope, regime class) is a pure
function of the mutation mask and is read from compile-time lookup tables with no
data-dependent branch in the hot loop; the whole working set ($\approx 4.4$~KiB)
is L1-resident. A struct-of-arrays batch interface processes ensembles of
trajectories in one call and is written to auto-vectorize, with an optional
explicit-SIMD path that reproduces the scalar output bit-for-bit. Trajectories
over the eight-symbol alphabet $Q_3$ form a four-level hierarchy whose
run-collapsed \emph{shape} (with an aligned \emph{dwell} vector) is the canonical
form, while coarser projections (net mutation, endpoint code) serve only as
pre-filters and never as form identities. The kernel performs no inference: states
are produced by an explicit, versioned anchoring map from observations, and a
language model may only synthesize policy or explanation over signatures the
kernel has already computed. We give the formal model with short proofs, the
computational architecture and its memory rationale, an exhaustive
cross-implementation parity methodology, the evaluation hypotheses, the limits,
and explicit non-claims (in particular: no complexity-theory claim, and no model
of quantum evolution).
\end{abstract}

\section*{Sommario esteso}
\addcontentsline{toc}{section}{Sommario esteso (Italiano)}
SemaGraph è un kernel deterministico che rappresenta il cambiamento osservato come
traiettorie finite su un alfabeto piccolo e fisso, e che classifica e comprime
tali traiettorie usando soltanto operazioni intere native. Ogni stato è una parola
di sei bit ($Q_6$, $64$ stati); ogni transizione diretta tra due stati è una
maschera XOR $m = s \oplus t$; la base completa delle transizioni è costituita
dalle $64 \times 64 = 4096$ coppie ordinate. I metadati di una transizione
(distanza di Hamming, distanze per-modulo, ambito di modulo, classe di regime)
sono funzione esclusiva della maschera e vengono letti da tabelle di lookup
costruite a tempo di compilazione, senza alcun salto condizionato sui dati nel
ciclo caldo; l'insieme di lavoro ($\approx 4.4$~KiB) risiede in cache L1.

La tesi è volutamente ristretta: la parte \emph{strutturale} del ragionamento
sulle traiettorie è riconducibile a primitive elementari della CPU (XOR, popcount,
lookup) e parallelizzabile su grandi insiemi (\emph{ensemble}) di traiettorie,
lasciando al modello linguistico soltanto il lavoro \emph{semantico} (spiegazione,
sintesi di policy). Il modello linguistico non calcola mai stati, maschere o
firme. Le traiettorie sull'alfabeto a otto simboli $Q_3$ formano una gerarchia a
quattro livelli la cui \emph{forma} (lo \emph{shape} ottenuto collassando le
ripetizioni adiacenti, con il vettore \emph{dwell} allineato) è l'identità di
forma; le proiezioni più grossolane (mutazione netta, codice di endpoint) sono
solo pre-filtri e non possono mai fungere da chiave di forma. Lo stato non viene
inferito linguisticamente: è prodotto da una mappa di ancoraggio esplicita e
versionata a partire dalle osservazioni.

Il documento fornisce il modello formale con brevi dimostrazioni, l'architettura
computazionale e la sua motivazione rispetto alla gerarchia di memoria, la
metodologia di verifica di parità esaustiva fra l'implementazione di riferimento
in TypeScript e il core in Rust, le ipotesi di validazione, i limiti e le
non-affermazioni esplicite. In particolare SemaGraph \textbf{non} è un motore
fisico, non è un linguaggio di programmazione generale, non sostituisce solutori
numerici o statistici, non avanza alcuna tesi di teoria della complessità (nessuna
affermazione su P contro NP) e \textbf{non} modella l'evoluzione quantistica:
quando applicato a hardware quantistico, opera unicamente sulle sequenze classiche
di esiti di misura, cioè dopo il collasso.

\tableofcontents
\bigskip

\section{Introduction and motivation}
Many complex systems produce \emph{sequences} of states rather than isolated
states. Organizational processes, financial models, discretized physical events
and stochastic simulations can be observed as trajectories: an initial
configuration changes, passes through intermediate regimes, and reaches a final
configuration. The question is not only how to classify the current state, but how
to represent the way the system got there.

The concrete problem that motivated SemaGraph was sampling cost. A stochastic
simulation kernel required, for statistical robustness, more rollouts than
consumer hardware could produce in reasonable time. The usual response --- run
more samples --- attacks the statistics but not the use one makes of the samples
already available. SemaGraph proposes a change of representation instead: collapse
each trajectory to a canonical form and measure the distribution over \emph{forms},
which stabilizes with far fewer rollouts than the distribution over raw paths,
because the form space is coarser while still preserving the order of regimes
visited.

The state alphabet is derived from the I~Ching hexagram system, used solely
because it is the oldest known combinatorial system that is natively binary: six
positions, two values each, with a modular $3+3$ structure. The cosmology is
discarded; only the combinatorial skeleton is kept. What remains is $Q_6$: $64$
states, $4096$ transitions, a complete basis for a six-dimensional Boolean space.

The thesis is deliberately narrow and is stated as such throughout: the
\emph{structural} part of trajectory reasoning reduces to native CPU primitives
and parallelizes over ensembles, leaving only \emph{semantic} work to any
downstream language model.

\section{Formal model}
Notation: $B=\{0,1\}$, $\oplus$ is bitwise XOR, $\operatorname{popcount}(x)$ is the
Hamming weight of $x$, and $[a,b)$ is a half-open integer interval.

\subsection{The composite state space \texorpdfstring{$Q_6$}{Q6}}
\begin{equation}
Q_6 = B^6, \qquad |Q_6| = 2^6 = 64.
\end{equation}
A state $s\in Q_6$ is a six-bit word, stored as the low six bits of a byte;
the canonical string form is \texttt{S64-} followed by six binary digits. Each
state splits into two three-bit modules, written MSB-first as $b_5b_4b_3b_2b_1b_0$:
\begin{equation}
L(s) = (s \gg 3)\ \&\ 7, \qquad U(s) = s\ \&\ 7 .
\end{equation}
The split carries no inherited meaning; it is the combinatorial property that lets
a transition be classified as internal to one module or crossing both.

\subsection{Transitions, distances, scope}
For $s,t\in Q_6$ the \emph{mutation mask} and \emph{Hamming distance} are
\begin{equation}
m(s,t) = s\oplus t, \qquad d(s,t) = \operatorname{popcount}(s\oplus t)\in\{0,\dots,6\},
\end{equation}
with per-module distances $d_L = \operatorname{popcount}(L(s)\oplus L(t))$ and
$d_U = \operatorname{popcount}(U(s)\oplus U(t))$, and $d = d_L + d_U$. The
\emph{module scope} is \texttt{none}, \texttt{lower\_only}, \texttt{upper\_only} or
\texttt{both\_modules} according to which of $d_L,d_U$ are positive.

\subsection{Regime taxonomy}
Let $\mathit{both} = (d_L>0 \wedge d_U>0)$. The regime class is the first matching
clause, in this order:
\begin{center}
\begin{tabular}{ll}
\toprule
condition & class (code) \\
\midrule
$d=0$ & \texttt{no\_change} (0) \\
$d=1$ & \texttt{bit\_adjustment} (1) \\
$d\le 2 \wedge \neg\mathit{both}$ & \texttt{module\_reconfiguration} (2) \\
$d=6$ & \texttt{full\_bit\_reversal} (5) \\
$d\ge 4$ & \texttt{near\_total\_inversion} (4) \\
otherwise & \texttt{cross\_module\_regime\_shift} (3) \\
\bottomrule
\end{tabular}
\end{center}
Thus $d=2$ confined to one module is \texttt{module\_reconfiguration}, while $d=2$
split across both modules, and every $d=3$, is
\texttt{cross\_module\_regime\_shift}. The taxonomy is a total function of
$(d,d_L,d_U)$; the integer codes are part of the cross-implementation contract.

\subsection{Index and packed word}
The $64\times64=4096$ ordered pairs form the complete direct basis (the $64$
identities included; $64\times63=4032$ are non-identity). The canonical index is
\begin{equation}
I(s,t) = 64\,s + t \in [0,4096).
\end{equation}
For low-overhead FFI the full record is packed into one 64-bit word, least
significant field first:
\begin{lstlisting}
source:6 | target:6 | mutation:6 | distance:3 | d_L:2 | d_U:2 | scope:2 | regime:3 | index:12
\end{lstlisting}
This word is verified bit-for-bit against the reference for all $4096$ transitions.

\subsection{Chains and their compression}
A chain is an ordered sequence $C=(s_0,\dots,s_n)$, each step yielding
$m_i = s_{i-1}\oplus s_i$ and the ordered signature $(m_1,\dots,m_n)$.

\begin{proposition}[Net mutation telescopes]
$M_{\mathrm{net}} := m_1\oplus\cdots\oplus m_n = s_0 \oplus s_n.$
\end{proposition}
\begin{proof}
$\bigoplus_{i=1}^{n}(s_{i-1}\oplus s_i) = s_0\oplus s_n$, since every interior
$s_i$ occurs exactly twice and cancels under XOR. \end{proof}

Hence $M_{\mathrm{net}}$ depends only on the endpoints and discards order: distinct
chains can share a net mutation, which is precisely why the ordered signature and
the net mutation are retained as distinct objects. The \emph{cumulative distance}
$D_C = \sum_i \operatorname{popcount}(m_i)$ measures total traversal effort (it
does not cancel out-and-back moves). The operative signature is
$\Sigma(C) = (s_0, s_n, (m_1,\dots,m_n), M_{\mathrm{net}}, D_C)$; it indexes the raw
chain but never replaces it.

\subsection{The four-level trajectory hierarchy over \texorpdfstring{$Q_3$}{Q3}}
Trajectories are typed over the eight-symbol alphabet $Q_3=\{0,\dots,7\}$.

\textbf{Level 0 --- alphabet.} The eight elementary states.

\textbf{Level 1 --- oriented edge $T_{64}$.} The ordered pair $e(a,b)=8a+b\in[0,64)$.
\begin{proposition}[Recoverability]
$e$ is injective and orientation-preserving: $a=\lfloor e/8\rfloor$, $b=e\bmod 8$,
and $e(a,b)=e(b,a)$ iff $a=b$.
\end{proposition}
$T_{64}$ shares the cardinality $64$ with a $Q_6$ \emph{point} but is an oriented
$Q_3\times Q_3$ \emph{edge}; the two are distinct types and never interconvert.

\textbf{Level 2 --- paths, run-collapse, dwell.} A path is a word over $Q_3$. Its
run-collapsed \emph{shape} merges adjacent duplicate states; non-adjacent repeats
are preserved (so $[0,1,0,3]$ does not collapse and stays distinct from $[0,1,3]$).
The \emph{dwell} vector records, in run order, how many positions each run held.
\begin{proposition}[Lossless decomposition]
$\mathit{path}\mapsto(\mathit{shape},\mathit{dwell})$ is a bijection onto pairs in
which $\mathit{shape}$ has no two adjacent entries equal and $\mathit{dwell}$ is a
positive-integer vector of equal length.
\end{proposition}
\begin{proof}
This is run-length encoding restricted to maximal constant runs. The
no-adjacent-equal condition makes run boundaries unique, so repeating
$\mathit{shape}[k]$ exactly $\mathit{dwell}[k]$ times inverts the map. \end{proof}

\textbf{Level 3 --- declared lossy projections.} Each with its collapse factor:
\begin{center}
\begin{tabular}{lll}
\toprule
projection & image size & role \\
\midrule
net mutation ($Q_3$) & $8$ & coarse pre-filter only \\
endpoint code ($T_{64}$) & $64$ & medium pre-filter only \\
endpoint Hamming & scalar & magnitude only \\
\texttt{shape\_key} & injective on shape & \textbf{form identity} \\
\texttt{exact\_key} & injective on path & \textbf{trajectory identity} \\
\texttt{dwell\_signature} & --- & duration identity \\
\bottomrule
\end{tabular}
\end{center}

\subsection{Invariants and the comparison cascade}
\begin{proposition}[Endpoint soundness]
Equal run-collapsed shape implies equal endpoint.
\end{proposition}
\begin{proof}
Run-collapse never removes the extreme positions, so the shape's first and last
entries are the path's first and last states, and the endpoint is a function of
those two. \end{proof}
The converse fails: $[0,1,3]$ and $[0,2,3]$ share the endpoint $e(0,3)$ but differ
in shape. Hence the endpoint is a \emph{sound pre-filter} (a mismatch proves the
shapes differ) but \textbf{never} a form decision. The binding invariant is that
\emph{no} lossy projection is a form key: \texttt{shape\_key} decides form. The
comparison of two trajectories returns exactly one relation:
\begin{center}
\begin{tabular}{ll}
\toprule
result & condition \\
\midrule
\texttt{different\_trajectory} & endpoints differ (short-circuit) \\
\texttt{different\_pattern} & else shapes differ \\
\texttt{equivalent} & else dwell signatures equal \\
\texttt{same\_form\_different\_duration} & otherwise \\
\bottomrule
\end{tabular}
\end{center}

\section{Deterministic anchoring}
States are not inferred by a model. They are produced by an explicit anchoring map
under a versioned rule set $R=(r_1,\dots,r_6)$:
\begin{equation}
A_R : X \to Q_6, \qquad b_i = r_i(x)\in\{0,1\},
\end{equation}
where $X$ is the space of observed/parametric inputs and each $r_i$ is a declared
threshold, Boolean condition or predicate. Given the same $x$ and the same $R$,
$A_R(x)$ is identical (replayability). Measurement uncertainty is recorded at the
measurement layer --- tolerance, unit, source, timestamp, rule version --- and does
not enter the transition kernel. A language model may explain a signature, propose
policy, or compare approved policies; it must never fabricate observations or alter
states, masks or signatures. This is the anti-inference boundary.

\section{Computational architecture}
\subsection{Branchless, cache-resident classification}
Every scalar the classifier needs is a function of the six-bit mask $s\oplus t$,
so the decision surface fits in a few compile-time tables indexed by that mask. No
\texttt{if}/\texttt{match} on the data appears in the hot loop. The working set is
\begin{center}
\begin{tabular}{lrr}
\toprule
table & entries & bytes \\
\midrule
distance / lower / upper / scope / regime (by mask) & $5\times 64$ & $320$ \\
regime by transition $[u8;4096]$ & $4096$ & $4096$ \\
\midrule
total & & $4416$ \\
\bottomrule
\end{tabular}
\end{center}
At $\approx 4.4$~KiB the tables occupy roughly $14\%$ of a typical $32$~KiB L1 data
cache, so once warm the common path is a few L1-latency loads plus an XOR and a
multiply-add, with no branch to mispredict. The tables are built from the same
\texttt{const fn} classifiers as the scalar path, so they cannot drift; an
exhaustive test asserts equality on all $4096$ transitions.

\subsection{Struct-of-arrays batch and SIMD}
The batch interface takes parallel input arrays (sources, targets) and writes into
caller-provided output columns, as a flat counted loop with no early return, no
per-iteration allocation and no data branch. This is the layout an auto-vectorizer
can widen. An optional, non-default \texttt{simd} feature provides an explicit SSE2
path that produces output identical to the scalar batch (asserted by test); the
scalar version is the source of truth. The trajectory batch compresses an ensemble
in one call into preallocated columns.

\subsection{Determinism, substrate, and interface}
The classification path uses no floating point and no randomness. The kernel is
implemented in dependency-free Rust with a C ABI (and is WASM-compatible); the
TypeScript implementation remains the executable specification. A single
zero-dependency CLI, \texttt{semagraph}, exposes the operations; its \texttt{--json}
output is a stable contract designed for consumption by automated agents (data on
stdout, errors on stderr, conventional exit codes; NDJSON streaming for batches).

\section{Verification methodology}
Parity is exhaustive, not sampled. The TypeScript reference generates a golden
fixture describing all $4096$ transitions with canonical integer codes and the
packed FFI word; a dependency-free Rust test asserts the core reproduces every
entry exactly. A second, seeded fixture covers Q3 trajectory projections (shape,
dwell, net mutation, endpoint, endpoint Hamming, cumulative distance), and the
branchless path is tested against the validated path on every transition.
Continuous integration regenerates both fixtures from the reference and fails on
any drift, so the two implementations cannot diverge silently.

\section{Performance: architectural expectations}
We state expectations, not measurements. A Criterion benchmark harness compares the
naive per-item path, the branchless table path, the scalar batch, the optional SIMD
batch, and end-to-end trajectory compression over a seeded ensemble, reporting
transitions/second and trajectories/second. The branchless path is expected to be
at least as fast as the naive path (constant-time L1 loads replacing a comparison
chain); the batch is expected to meet or exceed the per-item path (vectorizable
shape); the SIMD gain is bounded by the scalar table gather and is a demonstration,
not a claim. All four paths produce identical output, which is tested rather than
assumed. Concrete numbers belong in the repository's \texttt{BENCHMARKS.md},
recorded with CPU, compiler flags and seed; they are deliberately not invented
here.

\section{Relationship to known formalisms}
\begin{longtable}{p{3.4cm}p{5cm}p{5cm}}
\toprule
Family & Similarity & Difference \\
\midrule
Finite automata & Finite states and transitions & Oriented toward compressing
\emph{observed} trajectories, not recognizing languages \\
Statecharts & Guarded transitions & Does not orchestrate application/UI behaviour \\
Timed automata & Sequences with durations & Dwell is a discrete count, not a
real-valued clock; no timing constraints \\
Rule engines & Predicates and rules & Produces no general inference; it classifies
and compresses \\
Policy-as-code & Formalizable operational rules & Policy is strictly downstream of
the computed signature \\
Markov-switching / Monte Carlo & Trajectories and regimes & Does not sample paths;
it deterministically classifies given paths \\
\bottomrule
\end{longtable}

\section{Domains and evaluation hypotheses}
The framing is domain-neutral: a deterministic compression layer for
state-transition trajectories where sampling cost is the bottleneck and
path-dependence is structural. Worked translations accompany the repository for
(i)~stochastic simulation (grouping rollouts by \texttt{shape\_key} to stabilize a
histogram over forms); (ii)~regime-switching Monte Carlo and path-dependent
instruments, where in non-ergodic settings the ensemble mean does not equal the
time average and the distribution over typical trajectories is what matters
\cite{peters2019,hamilton1989}; (iii)~classical analysis of measurement-shot
sequences from quantum hardware (see the non-claim in \S\ref{sec:not});
(iv)~decision-making under deep uncertainty, where the anchoring map is the
quantization of qualitative parameters. The testable hypotheses are: shape/signature
processing reduces computation versus full-trajectory processing; signature-keyed
policy reuse improves consistency; the information lost to compression is bounded
and its false-equivalence cases are characterizable; and the set of domains
admitting direct deterministic anchoring (versus those requiring an uncertainty
envelope first) can be delineated.

\section{Limits and threats to validity}
The system's quality depends on the anchoring rules: a poor $A_R$ yields
meaningless states regardless of the kernel. Compression is lossy, and the most
critical risk is \emph{false equivalence} --- distinct chains sharing a net mutation
or even an endpoint. This is mitigated structurally (the form decision uses
\texttt{shape\_key}, not a lossy projection) and operationally (the raw chain is
always retained and the signature only indexes it), but it cannot be eliminated.
The kernel does not replace numerical or scientific solvers.

\section{What this is \emph{not}}\label{sec:not}
SemaGraph is not a physics engine, a general-purpose programming language, a
universal simulator, or a replacement for statistical/numerical models. It advances
\emph{no} complexity-theory result and makes no claim about P versus NP; it is an
engineering reduction of a narrow, structural task to native CPU primitives. It
does not model how the state axes are learned --- they are stipulated by the
anchoring map, not emergent. When applied to quantum hardware it does \textbf{not}
model quantum evolution (continuous unitary dynamics, complex amplitudes): it
operates only on the classical sequences of measurement outcomes after collapse,
as a classical signal-analysis layer.

\section{Future directions (speculative)}
Wider portable-SIMD and multi-lane implementations of the batch path; an
FPGA/lookup-ROM realization of the classifier as a near single-cycle lookup over the
12-bit transition index; and learned anchoring --- inferring the state axes from
data rather than stipulating them --- which is the genuinely hard, deliberately
out-of-scope problem.

\section{Conclusion}
SemaGraph formalizes a deterministic kernel that compresses observed change into
reproducible structural signatures over a finite $64$-state Boolean space. Its
contribution is architectural honesty about a narrow claim: the structural part of
trajectory reasoning is branchless, cache-resident and vectorizable, and can be
separated cleanly from the semantic work left to a language model. The mathematics,
the memory rationale, and an exhaustive parity methodology are given so that
mathematicians, engineers and physicists can each evaluate the part that concerns
them.

\begin{thebibliography}{9}
\bibitem{shannon1948} C.~E.~Shannon, ``A Mathematical Theory of Communication'',
\emph{Bell System Technical Journal}, 1948.
\bibitem{hamming1950} R.~W.~Hamming, ``Error Detecting and Error Correcting Codes'',
\emph{Bell System Technical Journal}, 1950.
\bibitem{hopcroft2006} J.~E.~Hopcroft, R.~Motwani, J.~D.~Ullman, \emph{Introduction
to Automata Theory, Languages, and Computation}, Pearson, 2006.
\bibitem{harel1987} D.~Harel, ``Statecharts: A Visual Formalism for Complex
Systems'', \emph{Science of Computer Programming}, 1987.
\bibitem{alur1994} R.~Alur, D.~L.~Dill, ``A Theory of Timed Automata'',
\emph{Theoretical Computer Science}, 1994.
\bibitem{drepper2007} U.~Drepper, ``What Every Programmer Should Know About
Memory'', technical report, 2007.
\bibitem{peters2019} O.~Peters, ``The ergodicity problem in economics'',
\emph{Nature Physics}, 2019.
\bibitem{hamilton1989} J.~D.~Hamilton, ``A New Approach to the Economic Analysis of
Nonstationary Time Series and the Business Cycle'', \emph{Econometrica}, 1989.
\bibitem{russell2021} S.~Russell, P.~Norvig, \emph{Artificial Intelligence: A Modern
Approach}, Pearson, 2021.
\end{thebibliography}

\end{document}
